\documentclass[11pt]{article}

\usepackage[english,italian]{babel}
\usepackage[a4paper, top=2cm, bottom=1.5cm, left=2cm, right=2cm]{geometry}
\usepackage{float}
\usepackage{ltablex}
\usepackage{titling}
\usepackage{blindtext}
\usepackage[utf8]{inputenc}
\usepackage[T1]{fontenc}
\usepackage{xcolor}
\usepackage{graphicx}
\usepackage{geometry}
\usepackage[italian]{babel}
\usepackage{tabularx}
\usepackage{longtable}
\usepackage{hyperref}
\usepackage[bottom]{footmisc}
\usepackage{fancyhdr}
\usepackage{titlesec}
\setcounter{secnumdepth}{4}
\usepackage{amsmath, amssymb}
\usepackage{array}
\usepackage{graphicx}
\usepackage{layouts}
\usepackage{url}
\usepackage{comment}
\usepackage{eurosym}






\begin{document}

	\thispagestyle{empty}
	\begin{titlepage}
		\begin{center}
			\includegraphics[scale = 0.05]{logo_unipd.png}\\
			\large \textbf{Università degli Studi di Padova} \\
			\vfill
			\includegraphics[scale = 0.7]{logo_small.jpg}\\
			\large \textbf{Bug Pharma} \\
			\vfill
			\large
			E-mail: 
			\texttt{bugpharma10@gmail.com}
			\vfill
			\Huge \textbf{Verbale interno del 02-11-2021}\\
			
			\large
			%per verbali esterni
			%\textbf{Zucchetti S.p.A. - Login Warrior}\\
			
			%per verbali interni
			\vfill
			\textbf{Ordine del giorno:} \\
			Individuazione aspetti positivi e negativi capitolato C2
			
			
			\vfill
			
			
			\begin{tabular}{r|l}
				\textbf{Approvazione} &  -\\
				\textbf{Redazione} &  \parbox[t]{5cm}{Lorenzo Piran}\\
				\textbf{Verifica} &  -\\
				\textbf{Stato} & Redatto \\
				\textbf{Uso} & Interno
			\end{tabular}
			\vfill
			
		\end{center}
	\end{titlepage}

	\section{Informazioni generali}
	\subsection{Luogo e data dell'incontro}
	\begin{itemize}
		\item \textbf{Luogo}: videoconferenza Zoom;
		\item \textbf{Data}: 02-11-2021;
		\item \textbf{Ora di inizio}: 13:30;
		\item \textbf{Ora di fine}: 14:15.
	\end{itemize}
	
	\subsection{Presenze}
	\begin{itemize}
		\item \textbf{Totale presenze}: 3;
		\item \textbf{Presenti}:
		\begin{itemize}
			\item Lorenzo Piran (segretario);
			\item Michele Masetto;
			\item Silvia Giro.
		\end{itemize}
		\item \textbf{Assenti}:
		\begin{itemize}
			\item Nicla Faccioli;
			\item Sara Nanni;
			\item Andrea Salmaso;
			\item Nicholas Sertori.
		\end{itemize}
		%\item \textbf{Partecipanti esterni}:
		%\begin{itemize}
		%	\item Gregorio Piccoli (Zucchetti S.p.A.).
		%\end{itemize}
	\end{itemize}

	\newpage

	\section{Punti salienti}
		Sono emersi i seguenti pro e contro, già in parte discussi nell'incontro di Mercoledì 20 Ottobre 2021:
		\begin{itemize}
			\item PRO:
			\begin{itemize}
				\item Il progetto permetterebbe al gruppo di cimentarsi in una sfida impegnativa, ma spendibile in futuro
				in termini di conoscenze acquisite;
			\end{itemize}
			\item CONTRO:
				\begin{itemize}
				\item I componenti del gruppo hanno esperienza quasi nulla sugli argomenti chiave che dovranno essere
				trattati, fatto che potrebbe portare a diversi inconvenienti.
				\item Sono emerse le seguenti domande, le cui risposte potrebbero portare ad un numero eccessivamente
				alto di casi d'uso, di possibile difficile gestione:
				\begin{enumerate}
					\item Cosa succederebbe se ci fosse un pacco smarrito? (es: problema di logistica, anche se il venditore
					ha inviato il pacco e l'acquirente ha pagato - soldi bloccati nello smart contract);
					\item Serve anonimato? È possibile ricontattare il venditore, per natura stessa della blockchain?;
					\item Che politica abbiamo per resi e rimborsi? E come possiamo esser sicuri che il venditore una volta
					ricevuto il pagamento non scappi coi soldi?;
					\item Come confermare la ricezione? (Ad esempio, se si utilizza la scansione del QR code, cosa succede
					se all'arrivo del pacco l'acquirente non ha la possibilità di fare lo scan per qualche motivo);
					\item E se volessi delegare qualcuno per ritirare il pacco?;
					\item È fattibile scegliere questo capitolato senza avere conoscenze pregresse sulle tecnologie che
					sembrano essere un po' il punto centrale?
				\end{enumerate}
			\end{itemize}
		\end{itemize}
\end{document}